\begin{spacing}{1.5}
\section{Conclusions}
PixelLab implements a course-aligned digital image processing laboratory with a local desktop interface and an offline Android application. The registry contains 59 operations, the course catalog contains 86 topics, and the repository retains standalone examples and original teaching references. The implementation joins validated controls, processed output, numerical analysis, plots and source examples in a repeatable workflow.

Fresh report checks exercised all default operation paths on a bounded sample and confirmed exact analysis-image recovery for four lossless codecs. The normalized equalization example reproduced the expected mappings. Retained logs provide additional evidence for desktop tests and Android integration. Together, these results support the software integration and numerical demonstration objectives. They do not establish measured learning gains or unrestricted device compatibility.

\section{Limitations}
The evaluation uses a supplied sample, exact classroom matrices, automated fixtures and retained emulator evidence. It has no large annotated image benchmark, participant study or physical ARM64 device validation. The fresh default execution checks cover one bounded input and do not exhaust the parameter space. Historical screenshots illustrate interface states and are not a fresh capture of every final feature.

Several experiments use reduced analysis images for practical execution. Lossless reconstruction applies to a bounded grayscale image, and selected segmentation results refer to resized coordinates. The course catalog has no explicit chapter-eight topic records despite morphology being executable through group mapping. Desktop and Android use different dependency versions, which creates a continuing need for numerical and regression checks on both environments.

\section{Future Recommendations}
A useful next step is a physical-device evaluation across several ARM64 phones, recording operation time, peak memory, failure behavior and output agreement with desktop. A reproducible numerical corpus should include constant, low-contrast, noisy, high-resolution and unusual-format images. Tests should retain input hashes, parameters, package versions and tolerances so that future changes can be assessed fairly.

The course catalog can be completed with explicit morphology topic records and reviewed page-level source mappings. Experiment history could persist parameters and input identifiers to support repeatable lab submissions. A separate student study could measure task completion and conceptual understanding with defined comparison conditions and consent. Production distribution would also benefit from release signing, accessibility review and a documented update process.

\end{spacing}